% Created 2026-05-11 Mon 02:22
% Intended LaTeX compiler: lualatex
\documentclass[11pt]{article}

\usepackage{amsmath}
\usepackage{fontspec}
\usepackage{graphicx}
\usepackage{longtable}
\usepackage{wrapfig}
\usepackage{rotating}
\usepackage[normalem]{ulem}
\usepackage{capt-of}
\usepackage{hyperref}
\usepackage{minted}
\usepackage{fontspec}
\setmainfont{UT Sans}[
BoldFont={UT Sans Bold},
UprightFont={UT Sans Regular},
FontFace={m}{n}{UT Sans Medium}
]
% Fallback fake italic using a slant
\usepackage{etoolbox}
\newcommand{\fakeitalic}[1]{{\addfontfeatures{FakeSlant=0.18}#1}}
\AtBeginDocument{\renewcommand{\emph}[1]{\fakeitalic{#1}}}
\usepackage{minted}
\usepackage[a4paper,margin=3cm]{geometry}
\usepackage{lastpage}
\usepackage{fancyhdr}
\usepackage{lastpage}
\pagestyle{fancy}                % set fancy style
\fancyhf{}                        % clear all headers and footers
\cfoot{\thepage/\pageref*{LastPage}}  % center footer
\renewcommand{\headrulewidth}{0pt}    % remove header line
\renewcommand{\footrulewidth}{0pt}    % remove footer line
% also override plain page style for first page
\fancypagestyle{plain}{%
\fancyhf{}%
\cfoot{\thepage/\pageref*{LastPage}}%
\renewcommand{\headrulewidth}{0pt}%
\renewcommand{\footrulewidth}{0pt}%
}
\usepackage{url}
\author{Elias-Robert BAJCSI}
\date{}
\title{Laborator 6\\\medskip
\large PLFP}
\hypersetup{
 pdfauthor={Elias-Robert BAJCSI},
 pdftitle={Laborator 6},
 pdfkeywords={},
 pdfsubject={},
 pdfcreator={},
 pdflang={English}}
\begin{document}

\maketitle
\section{\texttt{./Cargo.toml}}
\label{sec:orgdcec4de}
\begin{minted}[breaklines=true,breakanywhere=true,breaksymbol=,fontsize=\small,linenos=false]{toml}
[package]
name = "lab6"
version = "0.1.0"
edition = "2021"

[[bin]]
name = "ex1"
path = "src/bin/ex1.rs"
[[bin]]
name = "ex2"
path = "src/bin/ex2.rs"
[[bin]]
name = "ex3"
path = "src/bin/ex3.rs"
[[bin]]
name = "ex4"
path = "src/bin/ex4.rs"
\end{minted}
\section{\texttt{./src/bin/ex1.rs}}
\label{sec:orgf20b159}
\begin{minted}[breaklines=true,breakanywhere=true,breaksymbol=,fontsize=\small,linenos=false]{rs}
fn putere(baza: i64, exp: u32) -> i64 {
    if exp == 0 { 1 } else { baza * putere(baza, exp - 1) }
}

fn putere_tail(baza: i64, exp: u32, acc: i64) -> i64 {
    if exp == 0 { acc } else { putere_tail(baza, exp - 1, acc * baza) }
}

fn putere_tc(baza: i64, exp: u32) -> i64 { putere_tail(baza, exp, 1) }

fn inverseaza(s: &str) -> String {
    let chars: Vec<char> = s.chars().collect();
    fn aux(v: &[char]) -> String {
        if v.len() <= 1 { v.iter().collect() } else {
            let mut r = aux(&v[1..]);
            r.push(v[0]);
            r
        }
    }
    aux(&chars)
}

fn este_palindrom(s: &str) -> bool {
    let v: Vec<char> = s.chars().collect();
    fn aux(v: &[char]) -> bool {
        if v.len() <= 1 { true } else { v[0] == v[v.len() - 1] && aux(&v[1..v.len() - 1]) }
    }
    aux(&v)
}

fn medie_recursiva(sir: &[f64]) -> f64 {
    fn suma_rec(sir: &[f64]) -> f64 {
        if sir.is_empty() { 0.0 } else { sir[0] + suma_rec(&sir[1..]) }
    }
    if sir.is_empty() { 0.0 } else { suma_rec(sir) / sir.len() as f64 }
}

fn main() {
    println!("{} {}", putere(2, 10), putere(3, 0));
    println!("{} {}", putere_tc(2, 10), putere_tc(3, 0));
    println!("{}", inverseaza("recursiv"));
    println!("{} {}", este_palindrom("capac"), este_palindrom("rust"));
    println!("{}", medie_recursiva(&[1.0, 2.0, 3.0, 4.0]));
}
\end{minted}
\section{\texttt{./src/bin/ex2.rs}}
\label{sec:org321f9fe}
\begin{minted}[breaklines=true,breakanywhere=true,breaksymbol=,fontsize=\small,linenos=false]{rs}
#[derive(Debug)]
enum Arbore {
    Gol,
    Nod { val: i32, st: Box<Arbore>, dr: Box<Arbore> },
}

impl Arbore {
    fn contine(&self, val: i32) -> bool {
        match self {
            Arbore::Gol => false,
            Arbore::Nod { val: v, st, dr } => {
                if *v == val { true } else if val < *v { st.contine(val) } else { dr.contine(val) }
            }
        }
    }

    fn adancime(&self) -> usize {
        match self {
            Arbore::Gol => 0,
            Arbore::Nod { st, dr, .. } => 1 + st.adancime().max(dr.adancime()),
        }
    }
}

fn construieste_perechi<A: Clone, B: Clone>(a: &[A], b: &[B]) -> Vec<(A, B)> {
    if a.is_empty() || b.is_empty() {
        vec![]
    } else {
        let mut rest = construieste_perechi(&a[1..], &b[1..]);
        rest.insert(0, (a[0].clone(), b[0].clone()));
        rest
    }
}

fn main() {
    let arbore = Arbore::Nod {
        val: 8,
        st: Box::new(Arbore::Nod {
            val: 3,
            st: Box::new(Arbore::Nod { val: 1, st: Box::new(Arbore::Gol), dr: Box::new(Arbore::Gol) }),
            dr: Box::new(Arbore::Nod { val: 6, st: Box::new(Arbore::Gol), dr: Box::new(Arbore::Gol) }),
        }),
        dr: Box::new(Arbore::Nod {
            val: 10,
            st: Box::new(Arbore::Gol),
            dr: Box::new(Arbore::Nod { val: 14, st: Box::new(Arbore::Gol), dr: Box::new(Arbore::Gol) }),
        }),
    };

    println!("{} {}", arbore.contine(6), arbore.contine(9));
    println!("{}", arbore.adancime());
    println!("{:?}", construieste_perechi(&['a', 'b', 'c', 'd'], &[1, 2, 3, 4]));
}
\end{minted}
\section{\texttt{./src/bin/ex3.rs}}
\label{sec:org6ad6240}
\begin{minted}[breaklines=true,breakanywhere=true,breaksymbol=,fontsize=\small,linenos=false]{rs}
use std::collections::{HashMap, HashSet, VecDeque};

type Graf<'a> = HashMap<&'a str, Vec<&'a str>>;

fn bfs_rec<'a>(graf: &Graf<'a>, coada: &mut VecDeque<&'a str>, viz: &mut HashSet<&'a str>, out: &mut Vec<&'a str>) {
    if let Some(nod) = coada.pop_front() {
        out.push(nod);
        if let Some(vecini) = graf.get(nod) {
            for v in vecini {
                if viz.insert(*v) {
                    coada.push_back(*v);
                }
            }
        }
        bfs_rec(graf, coada, viz, out);
    }
}

fn bfs<'a>(graf: &Graf<'a>, start: &'a str) -> Vec<&'a str> {
    let mut coada = VecDeque::from([start]);
    let mut viz = HashSet::from([start]);
    let mut out = vec![];
    bfs_rec(graf, &mut coada, &mut viz, &mut out);
    out
}

fn toate_caile<'a>(graf: &Graf<'a>, start: &'a str, end: &'a str) -> Vec<Vec<&'a str>> {
    fn dfs<'a>(graf: &Graf<'a>, cur: &'a str, end: &'a str, viz: &mut HashSet<&'a str>, cale: &mut Vec<&'a str>, rez: &mut Vec<Vec<&'a str>>) {
        if cur == end {
            rez.push(cale.clone());
            return;
        }
        if let Some(vecini) = graf.get(cur) {
            for v in vecini {
                if viz.insert(*v) {
                    cale.push(*v);
                    dfs(graf, v, end, viz, cale, rez);
                    cale.pop();
                    viz.remove(v);
                }
            }
        }
    }

    let mut rez = vec![];
    let mut viz = HashSet::from([start]);
    let mut cale = vec![start];
    dfs(graf, start, end, &mut viz, &mut cale, &mut rez);
    rez
}

fn main() {
    let graf: Graf = HashMap::from([
        ("Ana", vec!["Bogdan", "Carmen"]),
        ("Bogdan", vec!["Dan", "Elena"]),
        ("Carmen", vec!["Elena"]),
        ("Dan", vec!["Filip"]),
        ("Elena", vec!["Filip"]),
        ("Filip", vec![]),
    ]);

    println!("{:?}", bfs(&graf, "Ana"));
    println!("{:?}", toate_caile(&graf, "Ana", "Filip"));
}
\end{minted}
\section{\texttt{./src/bin/ex4.rs}}
\label{sec:orgca51c9b}
\begin{minted}[breaklines=true,breakanywhere=true,breaksymbol=,fontsize=\small,linenos=false]{rs}
fn umple_rec(tabla: &mut [i32; 9], pas: usize, seed: &mut u64) {
    if pas == 9 {
        return;
    }
    *seed = seed.wrapping_mul(1664525).wrapping_add(1013904223);
    let mut idx = (*seed as usize) % 9;
    while tabla[idx] != -1 {
        idx = (idx + 1) % 9;
    }
    tabla[idx] = if pas % 2 == 0 { 1 } else { 0 };
    umple_rec(tabla, pas + 1, seed);
}

fn main() {
    let mut tabla = [-1; 9];
    let mut seed = 42u64;
    umple_rec(&mut tabla, 0, &mut seed);
    println!("{:?}", tabla);
}
\end{minted}
\end{document}
